\section{Literature Review and Conceptual Framework}

\subsection{Immigration Attitudes and Migrant Characteristics}

A large literature examines how host-country populations evaluate migrants and which characteristics shape 
these evaluations. One central distinction in this literature is based on whether attitudes toward migrants 
primarily reflect individuals' own economic interests and labor-market considerations or characteristics of the migrants.

Earlier economic studies provide some evidence that labor-market considerations are associated with 
immigration preferences \citep[e.g.,][]{scheve2001labor, mayda2006against}. Other studies, however, emphasize the 
importance of welfare concerns, cultural considerations, and group-related attitudes rather than just
individual labor-market interests \citep{burns2000economic, dustmann2007racial, sides2007european}. 
Reviewing this literature, \citet{hainmueller2014public} conclude that evidence for explanations based on 
individual economic self-interest is generally limited, particularly compared with broader sociotropic economic 
and cultural concerns.

Experimental evidence further challenges a simple labor-market competition explanation and rather points to the 
importance of migrants' own characteristics. \citet{hainmueller2010attitudes} find that both high- and low-skilled 
natives prefer high-skilled over low-skilled immigrants. If immigration attitudes were primarily driven by 
competition with similarly skilled workers, preferences would be expected to differ according to respondents' 
own skill levels. Instead, the common preference for high-skilled immigration suggests that respondents place 
considerable weight on migrants' human capital and potential economic contribution. This pattern is also found in 
multidimensional experiments. Higher education, occupational status, employability, and language proficiency are 
generally associated with more favorable evaluations of migrants \citep{hainmueller2015hidden, bansak2016economic,czymara2016deutschland}. 
More recent evidence from Germany similarly shows greater willingness to accept 
skilled labor migrants with recognized professional qualifications, good language skills, and a definite job offer 
\citep{wolff2025skilled}. Taken together, these studies show that characteristics associated with migrants' expected 
labor-market integration play an important role in how they are evaluated by the host-country population.

However, evaluations of migrants are not determined by human capital alone. Country of origin, cultural and religious 
background, and migration status also influence attitudes when other characteristics are held constant. 
\citet{hainmueller2015hidden} find differences in evaluations across immigrants' countries of origin alongside 
effects of education, occupation, and English proficiency. In the European context, \citet{bansak2016economic} show that 
preferences over asylum seekers depend jointly on economic, humanitarian, and cultural characteristics, with country 
of origin and religion affecting support in addition to employability and vulnerability. Similar patterns emerge in 
Germany. \citet{czymara2016deutschland} find that acceptance varies with migrants' country of origin and migration 
motive as well as their education and German-language proficiency. These findings suggest that respondents evaluate 
migrants based on both individual characteristics and the broader groups and migration categories to which they belong.
One possible explanation for the importance of such group-level characteristics is provided by theories of group 
position and perceived group threat. \citet{quillian1995prejudice}, building on the group-position perspective of \citet{blumer1958race}, 
argues that prejudice toward immigrant groups can be related to perceived collective threat and the relative position
of the in-group and out-group. In the European context, such perceptions are associated with factors including the 
relative size of immigrant groups and economic conditions. This perspective suggests that evaluations of migrants 
may partly reflect how migrant groups and their position within the host society are perceived.

Despite their differences, much of this literature shares a focus on attitudes and preferences toward migrants. 
Outcomes commonly concern support for immigration, willingness to admit or accept migrants, or preferences 
regarding their residence and employment. These outcomes capture how respondents evaluate migrants, but they 
do not necessarily reveal what respondents believe about migrants' likely outcomes. Attitudes and expectations are 
conceptually distinct. For example, a respondent may be willing to accept a refugee while expecting that the person 
will face difficulties finding employment or integrating socially. Conversely, a respondent may oppose further 
immigration while nevertheless expecting a particular migrant to succeed economically. Understanding how migrant 
characteristics affect beliefs about expected outcomes therefore requires examining these expectations directly, which
is what I aim to do in this thesis.

This thesis focuses on what outcomes respondents expect different migrants to experience, rather than on whether 
they are willing to accept them. I examine expectations concerning employment, occupational match, future income, 
discrimination, and social integration. The proposed factorial design makes it possible to examine whether these 
expectations vary with migrant characteristics, while also testing whether otherwise comparable migrants are 
evaluated differently. The next section develops this distinction further by reviewing the 
literature on beliefs, expectations, and misperceptions about migrants.

\subsection{Beliefs and Expectations about Migrants.}

\subsection{Differential Perceptions of Ukrainian and Syrian Migrants.}

\subsection{Contact, Personal Experience, Information, and Memory.}

